% ============================================================
%  كتاب الفلسفة الشامل - السنة الثالثة ثانوي
%  التصريف: XeLaTeX (مرتين لتحديث الفهرس والغلاف)
% ============================================================
\documentclass[12pt, twoside, openright]{book}

\usepackage{polyglossia}
\setmainlanguage[numerals=maghrib]{arabic}
\setotherlanguage{english}

\usepackage{fontspec}
\newfontfamily\arabicfont[Script=Arabic, Scale=1.2]{Amiri}
\newfontfamily\arabicfonttt[Script=Arabic, Scale=1.0]{Amiri}

\usepackage[
    paperwidth=17cm, paperheight=24cm,
    top=2.5cm, bottom=2.5cm,
    inner=2.5cm, outer=2.0cm,
    headheight=15pt
]{geometry}

\usepackage{xcolor}
\usepackage{setspace}
\usepackage{enumitem}
\usepackage{tikz}
\usetikzlibrary{calc}
\usepackage[most]{tcolorbox}
\setstretch{1.3}

% ---------- الألوان: لون لكل محور ----------
\definecolor{navy}{HTML}{1F2A44}
\definecolor{colA}{HTML}{1B8A8F}   % الأسرة والعمل والقيم والعدالة
\definecolor{colB}{HTML}{9B2D4A}   % الدولة والديمقراطية والسياسة والأخلاق
\definecolor{colC}{HTML}{C97A12}   % العلوم الإنسانية والتاريخ
\definecolor{colD}{HTML}{6B4C9A}   % اللغة والفكر والمنطق والرياضيات والحقيقة
\colorlet{chapcol}{navy}
\newcommand{\setchap}[1]{\colorlet{chapcol}{#1}}

% ---------- الرؤوس والتذييلات ----------
\usepackage{fancyhdr}
\pagestyle{fancy}
\fancyhf{}
\fancyhead[LE,RO]{\thepage}
\fancyhead[LO]{\leftmark}
\fancyhead[RE]{\rightmark}
\renewcommand{\headrule}{{\color{chapcol}\hrule width\headwidth height 0.8pt}}
\renewcommand{\chaptermark}[1]{\markboth{#1}{}}
\renewcommand{\sectionmark}[1]{\markright{#1}}
\fancypagestyle{plain}{\fancyhf{}\fancyfoot[C]{\color{chapcol}\thepage}\renewcommand{\headrulewidth}{0pt}}

% ---------- عنوان المحور ----------
\addto\captionsarabic{\renewcommand{\chaptername}{المحور}}
\usepackage{titlesec}
\newcommand{\chapterbanner}[1]{%
  \begin{tcolorbox}[colback=chapcol, colframe=chapcol, arc=3mm,
      halign=center, top=14pt, bottom=14pt]
    \color{white}
    \ifnum\value{chapter}>0
      {\large \chaptername\ \thechapter}\par\vspace{6pt}
    \fi
    {\Huge\bfseries #1}
  \end{tcolorbox}}
\titleformat{\chapter}[block]{\normalfont}{}{0pt}{\chapterbanner}
\titlespacing*{\chapter}{0pt}{0pt}{25pt}
\titleformat{\section}{\normalfont\Large\bfseries\color{chapcol}}{\thesection}{0.8em}{}
\titleformat{\subsection}{\normalfont\large\bfseries\color{chapcol!80!black}}{}{0pt}{}

% ---------- الصناديق ----------
\newtcolorbox{problembox}{enhanced, breakable, colback=chapcol!6,
  colframe=chapcol, fonttitle=\bfseries, title={الإشكالية},
  arc=2mm, boxrule=0.8pt}
\newtcolorbox{positionbox}[1]{enhanced, breakable, colback=white,
  colframe=chapcol!70, coltitle=white, colbacktitle=chapcol!70,
  fonttitle=\bfseries, title={#1}, arc=2mm, boxrule=0.6pt}
\newtcolorbox{synthbox}{enhanced, breakable, colback=chapcol!12,
  colframe=chapcol, fonttitle=\bfseries, title={التركيب وحل المشكلة},
  arc=2mm, boxrule=1pt}
\newtcolorbox{defbox}[1]{enhanced, breakable, colback=chapcol!6,
  colframe=chapcol, fonttitle=\bfseries, title={#1},
  arc=2mm, boxrule=0.8pt}

\begin{document}

% ============================================================
%  الغلاف
% ============================================================
\begin{titlepage}
\thispagestyle{empty}
\begin{tikzpicture}[remember picture, overlay]
  \fill[navy] (current page.south west) rectangle (current page.north east);
  % دوائر زخرفية بألوان المحاور
  \fill[colA, opacity=0.85] ([xshift=1cm, yshift=-2cm]current page.north west) circle (3.2cm);
  \fill[colB, opacity=0.80] ([xshift=-1cm, yshift=-4cm]current page.north east) circle (2.4cm);
  \fill[colC, opacity=0.80] ([xshift=2cm, yshift=-9cm]current page.north west) circle (1.5cm);
  \fill[colD, opacity=0.85] ([xshift=-2cm, yshift=-11cm]current page.north east) circle (1.9cm);
  % العنوان
  \node[text width=13cm, align=center, text=white]
    at ($(current page.center)+(0,3.2)$)
    {{\fontsize{44}{52}\selectfont\bfseries كتاب الفلسفة}\\[8pt]
     {\fontsize{30}{36}\selectfont الشامل}\\[14pt]
     {\Large موجه لتلاميذ السنة الثالثة ثانوي}};
  % المحاور الأربعة
  \foreach \i/\col/\txt in {
      1/colA/الأسرة والعمل والقيم والعدالة,
      2/colB/الدولة والديمقراطية والسياسة والأخلاق,
      3/colC/العلوم الإنسانية والتاريخ,
      4/colD/اللغة والفكر والمنطق والرياضيات والحقيقة}
    \node[fill=\col, text=white, minimum width=11.5cm, minimum height=1cm,
          rounded corners=4pt, font=\large]
      at ($(current page.center)+(0,{-1.35*\i+0.2})$) {\txt};
  % التذييل
  \node[text=white, align=center]
    at ($(current page.south)+(0,2.6)$)
    {{\large تأليف}\\[4pt]{\LARGE \textbf{الأستاذة وهيبة سعد السعود}}\\[8pt]
     {\small السنة الدراسية: 2026/2025}};
\end{tikzpicture}
\end{titlepage}

% ============================================================
%  المقدمات والمنهجية
% ============================================================
\frontmatter
\pagenumbering{arabic}

\chapter{تنبيه وإرشادات}
\begin{center}\textbf{تأليف الأستاذة وهيبة سعد السعود}\end{center}

رُتّب هذا الكتاب على أربعة محاور كبرى، لكل محور لونه الخاص: الأسرة والعمل والقيم والعدالة؛ الدولة والديمقراطية والسياسة والأخلاق؛ العلوم الإنسانية والتاريخ؛ اللغة والفكر والمنطق والرياضيات والحقيقة.

\begin{defbox}{طريقة عرض كل درس}
يبدأ كل درس بـ\textbf{الإشكالية}، ثم \textbf{المواقف الفلسفية} المتقابلة مع نقدها، وينتهي بـ\textbf{التركيب وحل المشكلة}.
\end{defbox}

\vspace{1em}
عناوين الدروس قابلة للتعديل لتوافق التدرج السنوي الرسمي وتدرج أستاذ المادة. والأمثلة والمواقف المذكورة مختصرة، ويُنصح بمراجعتها وإثرائها بالنصوص المقررة.

\tableofcontents

\chapter{فصل المنهجية الفلسفية}

\section{المقالة الجدلية}
\begin{defbox}{خطوات المقالة الجدلية}
\begin{enumerate}
  \item \textbf{المقدمة:} تمهيد، ثم طرح المشكلة، ثم صياغتها في سؤال.
  \item \textbf{العرض:} الأطروحة الأولى وحججها، ثم نقدها؛ ثم الأطروحة المقابلة وحججها، ثم نقدها؛ ثم التركيب.
  \item \textbf{الخاتمة:} حل المشكلة بموقف شخصي معلَّل.
\end{enumerate}
\end{defbox}

\section{تحليل النص الفلسفي}
\begin{defbox}{خطوات التحليل}
\begin{enumerate}
  \item التعريف بالنص وصاحبه، واستخراج المشكلة التي يعالجها.
  \item تحديد الأطروحة والحجج والأمثلة والمفاهيم المحورية.
  \item مناقشة النص: ما له وما عليه.
  \item خلاصة تبيّن قيمة النص وحدوده.
\end{enumerate}
\end{defbox}

\mainmatter

% ============================================================
%  المحور الأول
% ============================================================
\setchap{colA}
\chapter{الأسرة، العمل، القيم، والعدالة}

\section{الأسرة: طبيعة أم بناء اجتماعي؟}
\begin{problembox}
هل الأسرة مؤسسة طبيعية تفرضها الحاجة البيولوجية، أم بناء اجتماعي تصنعه الثقافة والقوانين؟
\end{problembox}
\subsection*{المواقف الفلسفية}
\begin{positionbox}{الأطروحة الأولى: الأسرة ظاهرة طبيعية}
عدّ أرسطو الأسرة أول جماعة تنشأ من اجتماع الرجل والمرأة لحفظ النوع، وتقوم على حاجة الطفل الطويلة إلى الرعاية. وربط ابن خلدون بين رابطة النسب وقوة العصبية. \textbf{النقد:} الطبيعة لا تفسّر تنوّع أشكال الأسرة بين المجتمعات والعصور.
\end{positionbox}
\begin{positionbox}{الأطروحة الثانية: الأسرة بناء اجتماعي}
بيّن دوركايم أن الأسرة مؤسسة تتغير بتغير المجتمع، من الأسرة الممتدة إلى النووية. ورأى ليفي شتراوس أن تحريم زواج الأقارب قاعدة ثقافية تنقل الإنسان من الطبيعة إلى الثقافة وتنشئ التحالف بين الجماعات. \textbf{النقد:} إهمال الأساس البيولوجي للتناسل والرعاية.
\end{positionbox}
\begin{synthbox}
للأسرة أساس بيولوجي هو التناسل ورعاية الصغار، لكن أشكالها وأدوارها وقواعدها تحددها الثقافة والمجتمع.
\end{synthbox}

\section{العمل: عقوبة أم قيمة؟}
\begin{problembox}
هل العمل عبء يحطّ من قيمة الإنسان، أم هو النشاط الذي يحقق إنسانيته؟
\end{problembox}
\subsection*{المواقف الفلسفية}
\begin{positionbox}{الأطروحة الأولى: العمل عبء}
احتقر التصور اليوناني القديم العمل اليدوي وتركه للعبيد، وخصّ الأحرار بالتأمل والسياسة. وبيّن ماركس أن العمل في النظام الرأسمالي يتحول إلى اغتراب: العامل منفصل عن منتجه وعن عمله وعن ذاته. \textbf{النقد:} هذا حكم على ظروف العمل لا على العمل ذاته.
\end{positionbox}
\begin{positionbox}{الأطروحة الثانية: العمل قيمة}
رأى هيغل أن العبد بعمله يحوّل الطبيعة ويكوّن وعيه بذاته فيتحرر. وعدّ ابن خلدون العمل أصل الكسب والقيمة وقيام العمران. \textbf{النقد:} لا يصح تمجيد العمل المفروض الذي يستغل صاحبه.
\end{positionbox}
\begin{synthbox}
معنى العمل مرهون بشروطه: العمل الحر العادل الأجر يحقق الإنسانية، والعمل المفروض المستغَل يغتربها.
\end{synthbox}

\section{العدالة: مساواة أم استحقاق؟}
\begin{problembox}
هل تعني العدالة المساواة المطلقة بين الأفراد، أم إعطاء كل ذي حق حقه؟
\end{problembox}
\subsection*{المواقف الفلسفية}
\begin{positionbox}{الأطروحة الأولى: العدالة مساواة}
يتساوى البشر في الكرامة وأمام القانون. وتطمح الاشتراكية عند ماركس إلى مجتمع يعطي كل فرد بقدر طاقته ويلبّي حاجاته. \textbf{النقد:} المساواة المطلقة تتجاهل اختلاف الجهد والكفاءة.
\end{positionbox}
\begin{positionbox}{الأطروحة الثانية: العدالة استحقاق}
ميّز أرسطو بين المساواة العددية والمساواة النسبية التي توزّع الخيرات بحسب الاستحقاق. وجعل رولز العدالة إنصافًا: حريات أساسية متساوية، وتفاوت مقبول بشرط أن يخدم الأقل حظًا وأن تتكافأ الفرص. \textbf{النقد:} معيار الاستحقاق نفسه موضع خلاف.
\end{positionbox}
\begin{synthbox}
العدالة مساواة في الحقوق وتكافؤ في الفرص، وتفاوت في المكافآت مبرَّر بالجهد ومحمي بدعم الضعفاء.
\end{synthbox}

% ============================================================
%  المحور الثاني
% ============================================================
\setchap{colB}
\chapter{الدولة، الديمقراطية، السياسة، والأخلاق}

\section{أصل الدولة وطبيعتها}
\begin{problembox}
هل الدولة نتاج اتفاق بين الأفراد، أم نتاج الغلبة والصراع بين الجماعات؟
\end{problembox}
\subsection*{المواقف الفلسفية}
\begin{positionbox}{الأطروحة الأولى: الدولة عقد اجتماعي}
رأى هوبز أن الأفراد تنازلوا عن حريتهم المطلقة لسلطة تحميهم من حرب الجميع ضد الجميع. وجعل لوك الدولة حارسة للحقوق الطبيعية، وجعلها روسو تعبيرًا عن الإرادة العامة. \textbf{النقد:} العقد افتراض نظري لا واقعة تاريخية.
\end{positionbox}
\begin{positionbox}{الأطروحة الثانية: الدولة غلبة وصراع}
ربط ابن خلدون قيام الدولة بالعصبية والغلبة، ثم تمر الدولة بأطوار الشباب فالهرم. وعدّها ماركس أداة تحمي مصالح الطبقة المسيطرة. \textbf{النقد:} تهمل وظائف الدولة في الأمن والعدل والخدمات.
\end{positionbox}
\begin{synthbox}
تنشأ الدولة غالبًا من ضرورة اجتماعية وظروف تاريخية، وتبقى شرعيتها رهينة بخدمة المصلحة العامة وصون الحقوق.
\end{synthbox}

\section{الديمقراطية: حكم الشعب أم خطر الأغلبية؟}
\begin{problembox}
هل الديمقراطية أفضل أنظمة الحكم، أم نظام تعتريه نقائص تهدد العدل والكفاءة؟
\end{problembox}
\subsection*{المواقف الفلسفية}
\begin{positionbox}{الأطروحة الأولى: نقد الديمقراطية}
انتقد أفلاطون حكم الجمهور لأنه يقدّم الأهواء على الكفاءة والمعرفة، فيفسح المجال للخطباء والمغامرين. \textbf{النقد:} الحكم بالمعرفة وحدها يفتح باب الاستبداد.
\end{positionbox}
\begin{positionbox}{الأطروحة الثانية: الدفاع عن الديمقراطية}
رأى روسو أن السيادة للشعب عبر الإرادة العامة. وبيّن بوبر أن قيمتها في تمكين المحكومين من تغيير حكامهم بلا عنف. \textbf{النقد:} قد تتحول سلطة الأغلبية إلى طغيان على الأقليات.
\end{positionbox}
\begin{synthbox}
الديمقراطية الأقدر على حفظ الحرية، بشرط أن تقترن بدولة القانون وحماية الأقليات وتعليم المواطنين.
\end{synthbox}

\section{السياسة والأخلاق}
\begin{problembox}
هل يجب أن تخضع السياسة للأخلاق، أم أن لها منطقها الخاص القائم على المصلحة والقوة؟
\end{problembox}
\subsection*{المواقف الفلسفية}
\begin{positionbox}{الأطروحة الأولى: استقلال السياسة عن الأخلاق}
فصل ميكيافيلي بين السياسة والأخلاق، فالأمير يحفظ دولته بما تقتضيه الظروف. \textbf{النقد:} يبرر الاستبداد ويفقد السلطة ثقة الناس.
\end{positionbox}
\begin{positionbox}{الأطروحة الثانية: خضوع السياسة للأخلاق}
رأى كانط أن السياسة الحقة لا تتقدم خطوة دون أن تؤدي الواجب الأخلاقي، وأن السلام الدائم قرين العدالة بين الدول. \textbf{النقد:} الأخلاق المجردة قد تعجز عن مواجهة تعقيدات الواقع.
\end{positionbox}
\begin{synthbox}
ميّز ماكس فيبر بين أخلاق القناعة وأخلاق المسؤولية؛ والسياسة الرشيدة تجمع بين المبادئ وحساب النتائج.
\end{synthbox}

% ============================================================
%  المحور الثالث
% ============================================================
\setchap{colC}
\chapter{العلوم الإنسانية والتاريخ}

\section{الظاهرة الإنسانية بين التفسير والفهم}
\begin{problembox}
هل يمكن دراسة الظاهرة الإنسانية دراسة علمية موضوعية كالظاهرة الطبيعية؟
\end{problembox}
\subsection*{المواقف الفلسفية}
\begin{positionbox}{الأطروحة الأولى: إمكان العلم الموضوعي}
دعا كونت إلى تطبيق المنهج الوضعي على المجتمع. وطالب دوركايم بدراسة الوقائع الاجتماعية كأشياء مستقلة عن الأفراد، واعتمد الإحصاء كما في دراسته للانتحار. \textbf{النقد:} الإنسان واعٍ وحر، لا يُختزل في شيء.
\end{positionbox}
\begin{positionbox}{الأطروحة الثانية: خصوصية الظاهرة الإنسانية}
ميّز دلتاي بين الطبيعة التي نفسّرها والحياة الإنسانية التي نفهمها. ومن العوائق: تعقّد الظاهرة وتغيرها، وتداخل ذات الباحث وموضوعه، وصعوبة التجريب لاعتبارات أخلاقية. \textbf{النقد:} لا يعني ذلك استحالة العلم بل تكييف مناهجه.
\end{positionbox}
\begin{synthbox}
العلوم الإنسانية علوم حقيقية بمنهج خاص يجمع التفسير والفهم، مع بذل الجهد لتقليل الذاتية.
\end{synthbox}

\section{التاريخ: هل هو علم؟}
\begin{problembox}
هل يملك التاريخ مقومات العلم، أم هو سرد ذاتي للماضي؟
\end{problembox}
\subsection*{المواقف الفلسفية}
\begin{positionbox}{الأطروحة الأولى: التاريخ ليس علمًا}
أحداث الماضي لا تتكرر ولا تُخضع للتجربة، ويُعرف الماضي بشهادات متفاوتة الصدق. والمؤرخ ابن عصره وعقيدته، وقد بيّن ريمون آرون حدود الموضوعية التاريخية. \textbf{النقد:} الصعوبة لا تلغي إمكان المنهج الصارم.
\end{positionbox}
\begin{positionbox}{الأطروحة الثانية: التاريخ علم}
انتقد ابن خلدون المؤرخين الذين ينقلون الأخبار دون تمحيص، ودعا إلى عرضها على طبائع العمران. ويعتمد المؤرخ نقد المصادر ومقارنة الوثائق وتعليل الأحداث. \textbf{النقد:} تبقى النتائج احتمالية لا قطعية.
\end{positionbox}
\begin{synthbox}
التاريخ علم له منهجه الصارم وحدوده، وموضوعيته نسبية يقاربها المؤرخ بالنقد والتحقق.
\end{synthbox}

\section{الحتمية والحرية في التاريخ}
\begin{problembox}
هل تسير الأحداث التاريخية وفق قوانين حتمية، أم يصنعها الإنسان بحريته؟
\end{problembox}
\subsection*{المواقف الفلسفية}
\begin{positionbox}{الأطروحة الأولى: الحتمية التاريخية}
جعل ماركس البنية الاقتصادية محركًا للتاريخ. ورأى ابن خلدون أن الدول تتعاقب وفق دورة محكومة بقوانين. ورأى هيغل أن العقل يتحقق في التاريخ بوسائل لا يدركها الأفراد. \textbf{النقد:} تلغي مسؤولية الإنسان وتهمل دور الأحداث الطارئة.
\end{positionbox}
\begin{positionbox}{الأطروحة الثانية: حرية الإنسان}
قال سارتر إن الإنسان هو ما يصنعه بنفسه، فالتاريخ حصيلة اختيارات حرة وقرارات حاسمة. \textbf{النقد:} الحرية مقيّدة بشروط مادية واجتماعية سابقة.
\end{positionbox}
\begin{synthbox}
يصنع الناس تاريخهم لكن ضمن شروط لم يختاروها؛ فالتاريخ تفاعل بين الضرورة والحرية.
\end{synthbox}

% ============================================================
%  المحور الرابع
% ============================================================
\setchap{colD}
\chapter{اللغة، الفكر، المنطق، الرياضيات، والحقيقة}

\section{العلاقة بين اللغة والفكر}
\begin{problembox}
هل اللغة شرط للتفكير، أم الفكر سابق عليها وهي مجرد أداة للتعبير؟
\end{problembox}
\subsection*{المواقف الفلسفية}
\begin{positionbox}{الأطروحة الأولى: اللغة تسبق الفكر وتشكّله}
عرّف سوسير العلامة اللغوية بأنها وحدة الدال والمدلول، ورأى أن الفكر قبل اللغة كتلة هلامية تنظمها الألفاظ. وذهبت فرضية سابير ووورف إلى أن لغة المرء تحدد طريقة إدراكه للعالم. \textbf{النقد:} نفكر أحيانًا دون كلمات، وتعجز الألفاظ عن بعض المعاني.
\end{positionbox}
\begin{positionbox}{الأطروحة الثانية: الفكر يسبق اللغة}
بيّن بياجيه أن ذكاء الطفل الحسي الحركي ينمو قبل الكلام. ويفكر الصم والبكم، وللحيوان نوع من الذكاء العملي. \textbf{النقد:} لا يبلغ الفكر المجرد درجته دون لغة.
\end{positionbox}
\begin{synthbox}
اللغة والفكر متداخلان: اللغة وعاء الفكر ووسيلة تجريده وتبليغه، والفكر أوسع من اللغة وهو ما يطوّرها.
\end{synthbox}

\section{المنطق الصوري وقيمته}
\begin{problembox}
ما قيمة المنطق الصوري في ضبط التفكير؟ وهل يكفي لبلوغ الحقيقة؟
\end{problembox}
\subsection*{المواقف الفلسفية}
\begin{positionbox}{الأطروحة الأولى: قيمة المنطق}
جعل أرسطو المنطق أداة للعلوم، وصاغ القياس ومبادئ التفكير: الهوية وعدم التناقض والثالث المرفوع. وهو يحمي العقل من المغالطة ويضمن صحة الاستدلال الصورية.
\end{positionbox}
\begin{positionbox}{الأطروحة الثانية: حدود المنطق}
صحة الاستدلال لا تعني صدق مقدماته، والقياس لا يضيف معرفة جديدة، وقد انتقده بيكون وديكارت لعقمه. \textbf{الرد:} طوّر بول وفريغه ورسل المنطق الرمزي الحديث.
\end{positionbox}
\begin{synthbox}
المنطق شرط ضروري للتفكير السليم غير كافٍ وحده؛ يحتاج إلى مادة صادقة من التجربة والعلم.
\end{synthbox}

\section{المعرفة الرياضية: عقل أم تجربة؟}
\begin{problembox}
هل الحقائق الرياضية قبلية مصدرها العقل وحده، أم مستمدة من التجربة؟
\end{problembox}
\subsection*{المواقف الفلسفية}
\begin{positionbox}{الأطروحة الأولى: الرياضيات عقلية}
اعتمد أفلاطون وديكارت على بداهة المبادئ الرياضية. وقال كانط إن الأحكام الرياضية تركيبية قبلية، أساسها صورتا الزمان والمكان في الذهن. \textbf{النقد:} ظهور الهندسات اللاإقليدية أثبت أن المسلمات ليست بديهيات مطلقة.
\end{positionbox}
\begin{positionbox}{الأطروحة الثانية: الرياضيات تجريبية واصطلاحية}
رأى ميل أن مبادئ الرياضيات نتاج استقراء التجربة. وبيّن بوانكاريه أن المسلمات اصطلاحات يختارها العقل لملاءمتها. \textbf{النقد:} الرياضيات لا تحتاج إلى التجربة لإثبات نتائجها.
\end{positionbox}
\begin{synthbox}
الرياضيات بناء استنتاجي فرضي يبدأ من مسلمات، وحقيقتها اتساقها الداخلي، وتصلح للتطبيق لأنها مجردة من الواقع.
\end{synthbox}

\section{معيار الحقيقة}
\begin{problembox}
ما المعيار الذي نميّز به الحقيقة من الوهم والخطأ؟
\end{problembox}
\subsection*{المواقف الفلسفية}
\begin{positionbox}{مطابقة الفكر للواقع}
الحقيقة عند أرسطو هي موافقة القول لما هو كائن، وتُختبر في العلوم التجريبية بالملاحظة والتجريب. \textbf{النقد:} الواقع لا يُدرك مباشرة بل عبر الحواس والنظريات.
\end{positionbox}
\begin{positionbox}{الاتساق الداخلي والبداهة}
اعتمد ديكارت على الوضوح والتميّز معيارًا، وتقوم الرياضيات على اتساق نسقها. \textbf{النقد:} الاتساق لا يضمن مطابقة الواقع.
\end{positionbox}
\begin{positionbox}{النفعية}
قال وليم جيمس إن الصادق ما ينجح في الحياة العملية. \textbf{النقد:} قد تنفع أوهام ولا تكون حقائق.
\end{positionbox}
\begin{synthbox}
يتعدد المعيار بتعدد المجالات: المطابقة والتجربة للعلوم الطبيعية، والاتساق للمنطق والرياضيات، والحقيقة نسبية في إدراكها، موضوعية في قصدها.
\end{synthbox}

\end{document}
